% ============================================================================
%  Preamble — LMU thesis template (English), adapted for a Master's thesis.
%  Merges the LMU layout packages with the mathematical packages our chapters
%  need. Loaded once from main.tex; document class (book, 12pt) is set there.
% ============================================================================

% ── LMU layout ───────────────────────────────────────────────────────────────
\usepackage{a4wide}
\usepackage{fancyhdr}
\usepackage{graphicx}
\graphicspath{{figures/}{./}}    % figures/ for plots, ./ for the seal
\usepackage[english]{babel}      % was: \usepackage{german}

% ── mathematics ──────────────────────────────────────────────────────────────
\usepackage{amsmath}
\usepackage{mathtools}
\usepackage{amssymb}
\usepackage{amsthm}
\usepackage{bm}

% ── tables, lists, floats ────────────────────────────────────────────────────
\usepackage{booktabs}
\usepackage{array}
\usepackage{enumitem}
\usepackage{caption}
\usepackage{xcolor}
\usepackage{tikz}
\usetikzlibrary{positioning}

% ── citations, then hyperref last, then todonotes ────────────────────────────
\usepackage[round]{natbib}
\usepackage[bookmarks,hidelinks]{hyperref}

% Status notes (scaffolding). For a clean PDF with no notes, change to
%   \usepackage[disable]{todonotes}
\usepackage[colorinlistoftodos,textsize=footnotesize,textwidth=2.6cm]{todonotes}
\newcommand{\towrite}[1]{\todo[color=yellow!70,inline]{\textbf{TO-WRITE:} #1}}
\newcommand{\stale}[1]{\todo[color=orange!60,inline]{\textbf{STALE:} #1}}
\newcommand{\port}[1]{\todo[color=green!35,inline]{\textbf{PORT:} #1}}

% ── running headers (from the LMU template) ──────────────────────────────────
\pagestyle{fancyplain}
\renewcommand{\chaptermark}[1]{\markboth{\thechapter.\ #1}{}}
\renewcommand{\sectionmark}[1]{\markright{\thesection\ #1}}
\lhead[\fancyplain{}{\bfseries\thepage}]{\fancyplain{}{\bfseries\rightmark}}
\rhead[\fancyplain{}{\bfseries\leftmark}]{\fancyplain{}{\bfseries\thepage}}
\cfoot{}

% ── theorem environments ─────────────────────────────────────────────────────
\theoremstyle{plain}
\newtheorem{theorem}{Theorem}[chapter]
\newtheorem{proposition}[theorem]{Proposition}
\newtheorem{lemma}[theorem]{Lemma}
\newtheorem{corollary}[theorem]{Corollary}
\theoremstyle{definition}
\newtheorem{assumption}{Assumption}[chapter]
\theoremstyle{remark}
\newtheorem{remark}{Remark}[chapter]
